% !TeX program = lualatex
% Compile twice: lualatex 107.tex
% All references and prose are embedded; no BibTeX database or external inputs.
\documentclass[11pt,a4paper]{article}
\usepackage[margin=24mm,headheight=14pt]{geometry}
\usepackage{fontspec}
\setmainfont{Latin Modern Roman}
\setsansfont{Latin Modern Sans}
\setmonofont{Latin Modern Mono}
\usepackage{amsmath,amssymb,mathtools}
\usepackage{unicode-math}
\setmathfont{Latin Modern Math}
\usepackage{microtype}
\usepackage{enumitem,xurl,needspace}
\usepackage[dvipsnames]{xcolor}
\usepackage{hyperref}
\hypersetup{unicode=true,colorlinks=true,linkcolor=MidnightBlue,urlcolor=MidnightBlue,
 pdftitle={Research attempt: Problem 107},pdfauthor={Research notebook},bookmarksnumbered=true}
\usepackage{bookmark}
\usepackage{tocloft}
\setlength{\cftsecnumwidth}{3em}
\cftsetpnumwidth{2.5em}
\cftsetrmarg{4em}
\usepackage{titlesec}
\titleformat{\section}{\Large\bfseries\raggedright}{\thesection}{0.7em}{}
\usepackage{fancyhdr}
\pagestyle{fancy}\fancyhf{}
\fancyhead[L]{\small\sffamily Open Problems Math | Research notebook}
\fancyhead[R]{\small Problem 107 | 27 September 2026}
\fancyfoot[C]{\thepage}
\setlength{\parindent}{0pt}
\setlength{\parskip}{3pt plus 1pt}
\setlength{\emergencystretch}{3em}
\setcounter{tocdepth}{1}
\setcounter{secnumdepth}{1}
\setlist[itemize]{leftmargin=3em,itemsep=5pt,topsep=4pt,parsep=0pt}
\allowdisplaybreaks
\newcommand{\N}{\mathbb{N}}
\newcommand{\Z}{\mathbb{Z}}
\newcommand{\Q}{\mathbb{Q}}
\newcommand{\R}{\mathbb{R}}
\newcommand{\C}{\mathbb{C}}
\newcommand{\F}{\mathbb{F}}
\newcommand{\A}{\mathbb{A}}
\newcommand{\PP}{\mathbb P}
\newcommand{\E}{\mathbb{E}}
\newcommand{\eps}{\varepsilon}
\newcommand{\dd}{\,\mathrm d}
\newcommand{\sref}[2]{\hyperlink{source:#1:#2}{[S#2]}}
\newcommand{\eref}[2]{\hyperlink{extra:#1:#2}{[E#2]}}
% Turn the next switch off for a shorter reading copy. The quotations preserve
% every exactTarget field, including qualifications omitted from a short summary.
\newif\ifshowcatalogue
\showcataloguetrue
\newcommand{\cataloguescope}[1]{\ifshowcatalogue
 \paragraph{Exact scope in the supplied catalogue.}
 \begin{quote}\small #1\end{quote}\fi}
\usepackage{etoolbox}
\AtBeginEnvironment{itemize}{\small}
\numberwithin{equation}{section}
% Standalone export. Original section 107 preserved verbatim outside marked additions.
\begin{document}
\setcounter{section}{106}
% ================================================================
% problemId: problem.unbounded-ranks-of-elliptic-curves-over-q
% source releaseRank: 107; source releaseStatus: open
% exactTarget SHA-256: 60fcb006eee42ee18f5f784fbd89b17f076d0f29a40d4750353a037cad3731e9
\clearpage
\section{\texorpdfstring{Unbounded Ranks of Elliptic Curves over Q}{Unbounded Ranks of Elliptic Curves over Q}}\label{107}
\subsection{Definitions and mathematical statement}
For an elliptic curve over $\Q$, write $E(\Q)\simeq\Z^{r(E)}\oplus T$ with $T$ finite. The target is
\[
 \forall r\in\Z_{\ge0}\ \exists E/\Q:\qquad r(E)\ge r,
\]
or $\sup_{E/\Q}r(E)=\infty$. The rank is that of the actual rational-point group. Analytic ranks, Selmer upper bounds, and ranks after enlarging the base field are not interchangeable with this integer. No algorithm producing the curves, or prescribed rate of growth in their heights, is required.
\par\smallskip\textit{Formulation and concept references:} \sref{107}{1}.
\cataloguescope{Prove that for EVERY nonnegative integer r there is SOME elliptic curve E/Q whose actual Mordell-Weil group E(Q) has Z-rank at least r. The base field is exactly Q and rank means the free abelian rank in Mordell's theorem, not analytic rank, a Selmer bound or random-matrix corank. Thus the set of actual ranks of ALL elliptic curves over Q is unbounded. Counting one representative per Q-isomorphism class does not alter this existence question. No effective construction, specified family, uniform curve-producing algorithm or predetermined rate of growth is required.}
\subsection{Short English statement}
Can elliptic curves over the rationals have arbitrarily many independent rational points of infinite order?
% BEGIN RESEARCH ADDITION 107
\subsection{Research attempt}

\paragraph{Current frontier.}
The source paper develops searches for high-rank curves with specified
isogenies; it does not establish unboundedness. \sref{107}{1}
The record table maintained by Dujella currently lists a 2026 example of
Alp\"oge--Howell with rank at least $31$, together with rational point data.
That is a reported finite rank certificate, not an assertion that the rank
is exactly $31$; its full independence computation has not been audited in
this notebook. \eref{107}{2}

There is also substantive evidence pointing the other way. The
Park--Poonen--Voight--Wood model predicts only finitely many curves of rank
greater than $21$, while explicitly warning that special arithmetic families
could depart from the model. This is a heuristic about actual curves, not
a theorem bounding their ranks. A finite example of rank $29$ or $31$ does
not contradict a prediction of \emph{finitely many} exceptional curves.
\eref{107}{1} The attempt below therefore tests both a construction mechanism
and the obstruction to descending it, without treating either expectation
as an established answer.

\paragraph{Attack route.}
One could attempt arbitrarily high generic rank over $\Q(t)$ and then
specialize, but the arbitrarily-high-generic-rank input would itself be a
major unresolved step. A second route creates points in independent quadratic
extensions and tries to descend them; this has an accessible construction
but an exact Galois obstruction. A third route searches for many points in
one square class and certifies their actual independence by descent maps.
I carry out the second route completely in its enlarged field, identify why
it does not answer the target, and then implement the third route over
$\Q$ with exact certificates.

\paragraph{Attempt.}
\textbf{1. Arbitrarily many independent points, but in the wrong field.}
Fix
\[
 E_0:\quad y^2=x^3-x,
 \qquad f(x)=x^3-x.
\]
For any prescribed $r\geq1$, choose distinct primes $q_1,\ldots,q_r\geq5$.
For each $i$, the Chinese remainder theorem provides an integer $x_i\geq2$
such that
\begin{equation}\label{107:crt}
 x_i\equiv q_i\pmod{q_i^2},\qquad
 x_i\equiv2\pmod{q_j}\quad(j\neq i).
\end{equation}
Put $d_i=f(x_i)>0$. Because $f'(0)=-1$ and, more directly,
$f(q_i)\equiv-q_i\pmod{q_i^2}$,
\[
 v_{q_i}(d_i)=1,\qquad v_{q_j}(d_i)=0\quad(j\neq i).
\]
The second assertion uses $f(2)=6$ and $q_j\geq5$. Thus the classes of
$d_1,\ldots,d_r$ are independent in $\Q^\times/\Q^{\times2}$: applying
$v_{q_j}$ modulo two to a square relation detects its $j$th exponent.
Consequently
\[
 K_r=\Q(\sqrt{d_1},\ldots,\sqrt{d_r})
\]
is multiquadratic of degree $2^r$, with an automorphism independently
changing each square-root sign. Let $P_i=(x_i,\sqrt{d_i})\in E_0(K_r)$.

\medskip
\textbf{2. Nontorsion without an unverified numerical height calculation.}
The following elementary height argument applies to each $P_i$. For a point
$P$ on $E_0$ whose $x$-coordinate is rational, the duplication formula is
\begin{equation}\label{107:double}
 x(2P)=\frac{(x(P)^2+1)^2}{4x(P)(x(P)^2-1)}.
\end{equation}
Write $x(P)=a/b$ in lowest terms and put $H(x)=\max(|a|,|b|)$. The numerator
and denominator before cancellation in \eqref{107:double} are
\[
 (a^2+b^2)^2,\qquad 4ab(a^2-b^2).
\]
Their greatest common divisor divides $4$. An odd common prime could divide
neither $a$ nor $b$, and would divide both $a^2+b^2$ and $a^2-b^2$, which is
impossible. If $a,b$ are both odd, the numerator has $2$-adic valuation
exactly $2$; if they have opposite parity, it is odd. Hence, for $H(x)\geq2$,
\begin{equation}\label{107:height}
 H(x(2P))\geq\frac{H(x(P))^4}{4}.
\end{equation}
The exceptional denominator values $x=0,1,-1$ have height one, so they cannot
interrupt this iteration. With $h=\log H$ it follows that
\[
 h(x(2^nP))\geq
 4^n\left(h(x(P))-\frac{\log4}{3}\right)+\frac{\log4}{3}.
\]
For $x(P)=x_i\geq2$, the coefficient of $4^n$ is positive. The iterated
$x$-coordinates therefore have unbounded height, proving that $P_i$ is
nontorsion. This argument is valid in the quadratic field of $P_i$: only
its $x$-coordinate, and those of its doubles, are required to be rational.

Now suppose $\sum_i n_iP_i=O$ with integers $n_i$. Apply the automorphism
that changes only the sign of $\sqrt{d_j}$ and subtract the original
relation. This gives $2n_jP_j=O$. Since $P_j$ is nontorsion, $n_j=0$.
Thus all $P_i$ are independent, and
\begin{equation}\label{107:wrongfield}
 \operatorname{rank}E_0(K_r)\geq r.
\end{equation}
Every assertion up to \eqref{107:wrongfield} has been proved; the field is
explicitly $K_r$, not $\Q$.

\medskip
\textbf{3. Why the apparent descent gives zero.}
For the sign character $\chi_i$ belonging to $d_i$, one has
$\sigma P_i=\chi_i(\sigma)P_i$. In particular
\[
 \operatorname{Tr}_{K_r/\Q}(P_i)
 =\sum_{\sigma\in\operatorname{Gal}(K_r/\Q)}\sigma P_i=O.
\]
The usual rational decomposition by characters also explains the failure:
after tensoring the point group with $\Q$, different characters give direct
summands, each identified with the rational point space on the corresponding
quadratic twist. The constructed points occupy \emph{different} summands.
One point on each of many twists is not many independent points on one
curve. Averaging the ranks over the $2^r$ characters gives no unbounded
lower bound for one summand.

To make this concrete, write $d_i=e_i z_i^2$ with $e_i$ squarefree.
The point $P_i$ corresponds to the rational point
\[
 (e_ix_i,e_i^2z_i)\quad\text{on}\quad
 Y^2=X^3-e_i^2X.
\]
The independent square classes force the $e_i$ to be different. Making all
$e_i$ equal would remove the very sign-changing automorphisms used to prove
independence. That is the exact missing bridge in this construction.

\medskip
\textbf{4. A same-twist search and a rank certificate using actual points.}
For a fixed squarefree $d>0$, put
\[
 E_d:\quad Y^2=X(X-d)(X+d).
\]
If $x\in\Q$ and $x^3-x=dz^2$ for $z\in\Q\setminus\{0\}$, then
\begin{equation}\label{107:point}
 (X,Y)=(dx,d^2z)\in E_d(\Q).
\end{equation}
For $x=a/b>1$ in lowest terms,
\[
 x^3-x=\frac{a(a-b)(a+b)}{b^3}
       =\frac{ab(a-b)(a+b)}{b^4}.
\]
Thus the squarefree part of $ab(a-b)(a+b)$ identifies the twist. This makes
it possible to group rational inputs by $d$ without mistaking an integer
coordinate bound for an exhaustive search over rational points.

Here is the exact independence test. For $P=(X,Y)$ not a point of order two,
set
\[
 \kappa(P)=([X],[X-d],[X+d])
 \in(\Q^\times/\Q^{\times2})^3,
\]
where brackets denote square classes, written multiplicatively. Put
$\kappa(O)=(1,1,1)$, and for a root $e\in\{0,d,-d\}$ replace the vanishing
coordinate of $\kappa((e,0))$ by
$\prod_{e'\neq e}(e-e')$; retain the nonvanishing coordinates. In particular,
\begin{equation}\label{107:torsion}
 \kappa(T_0)=([-1],[-d],[d]),\qquad
 \kappa(T_d)=([d],[2],[2d]),
\end{equation}
where $T_0=(0,0)$ and $T_d=(d,0)$.

\emph{Why this is a homomorphism.} A nonvertical rational line meeting the
cubic at $P,Q,R$ has $P+Q+R=O$. Substitute its equation into the cubic and
evaluate the resulting monic cubic polynomial at a root $e$. If none of
the three points is $(e,0)$, this gives
\[
 (X(P)-e)(X(Q)-e)(X(R)-e)=(\text{rational number})^2.
\]
If one point is $(e,0)$, differentiate that polynomial at $e$ instead. The
product of the other two $X-e$ factors is
$\prod_{e'\neq e}(e-e')$, which gives the same square-class identity using
the prescribed replacement. Tangencies at non-two-torsion points are handled
with intersection multiplicity. A vertical line pairs a point with its
negative and gives the identity because their $X$-coordinates agree; doubling
a point of order two is the same case. These cases cover the group law.
Thus $\kappa(P+Q)=\kappa(P)\kappa(Q)$, and $2E_d(\Q)$ is in its kernel.

Since $E_d(\Q)\simeq\Z^r\oplus T$ and $E_d[2](\Q)$ has dimension two over
$\mathbb F_2$,
\[
 \dim_{\mathbb F_2}\bigl(E_d(\Q)/2E_d(\Q)\bigr)=r+2.
\]
Consequently a matrix of square-class columns obtained from \emph{actual}
points, of binary rank $R$, proves
\begin{equation}\label{107:rankcert}
 \operatorname{rank}E_d(\Q)\geq R-2.
\end{equation}
No injectivity theorem for $\kappa$ is needed for this lower bound. It is not
a Selmer upper bound, an analytic-rank guess, or a conditional BSD argument.

\medskip
\textbf{5. A fully checkable rank-at-least-four example.}
The rational search $1\leq b<a\leq600$, $\gcd(a,b)=1$, finds the squarefree
integer
\[
 d=29274=2\cdot3\cdot7\cdot17\cdot41
\]
and the following four points on $E_d$:
\begin{center}
\begin{tabular}{c|r|r}
$x=X/d$ & $X$ & $Y$\\\hline
$24/17$ & 41328 & 5930568\\
$34/7$ & 142188 & 52467372\\
$41/34$ & 35301 & 3706605\\
$121/2$ & 1771077 & 2356659459
\end{tabular}
\end{center}
Call them $P_1,\ldots,P_4$ in this order. Direct substitution verifies their
curve equations. Representatives for their square-class triples are
\[
 (287,246,42),\quad (123,1394,102),\quad
 (21,123,287),\quad (14637,123,119).
\]
Use columns $T_0,T_d,P_1,P_2,P_3,P_4$. Select the following six binary
coordinates: the sign of the first square class, its valuations at
$2,3,7,17$, and the valuation at $17$ of the second square class. The
resulting matrix is
\begin{equation}\label{107:matrix}
 \begin{pmatrix}
 1&0&0&0&0&0\\
 0&1&0&0&0&0\\
 0&1&0&1&1&1\\
 0&1&1&0&1&1\\
 0&1&0&0&0&1\\
 1&0&0&1&0&0
 \end{pmatrix}.
\end{equation}
Its integer determinant is $-1$, so its rank over $\mathbb F_2$ is six.
Equation \eqref{107:rankcert} proves that this curve has actual rational
rank at least four. The two torsion columns are independent. Since $T/2T$ has dimension two,
they span the entire image of the rational torsion subgroup under $\kappa$.
The four point columns are independent modulo that span. In the free
quotient $E_d(\Q)/T$, divide a hypothetical nonzero integer relation by
the largest common power of two in its coefficients. Lifting back gives a
relation modulo torsion with an odd coefficient, contradicting those four
independent square-class columns. Thus the four points are independent
modulo torsion as well. No claim is made that the curve or
these points are new, or that its exact rank has been determined here.

The earlier integer-only search $2\leq x\leq30000$ produced certificates of
rank at least two, for example $d=210$ and $x=6,15$. Allowing rational inputs
was therefore a genuine enlargement of the experiment rather than a change
in the target. The companion script reproduces both searches and verifies
\eqref{107:matrix} using exact arithmetic.

\paragraph{Stress test.}
The enlarged-field construction proves its own nontorsion and independence
claims, but its trace calculation explicitly blocks the attempted descent.
The same-twist search fixes this field issue, at the cost of losing the
Galois-sign independence argument. Distinct points on a curve can be
multiples of one point; bucket size alone is therefore not a rank bound.
The matrix certificate, rather than the number of hits, is what proves the
finite rank statement.

The asymptotic missing statement in the restricted search is:
\begin{quote}
For every $r$ there is a squarefree $d>0$ and a finite collection of rational
solutions to $x^3-x=dz^2$ whose square-class columns, together with the
rational two-torsion columns, have binary rank at least $r+2$.
\end{quote}
This would imply the catalogue target, but it is stronger because it restricts
to the fixed-$j$ family $j=1728$ and imposes a particular certificate mechanism.
It is not known from the calculations. Even a proof that some ranks elsewhere
are unbounded would not automatically show that this particular restricted
family suffices. Conversely, failure of this route would not refute
unboundedness among \emph{all} elliptic curves.

The opposing boundedness model is not used as a theorem at any point.
Neither small average rank nor finitely many computational records decides
whether arbitrarily sparse high-rank exceptional curves exist. The model's
caution about special families is relevant precisely to that distinction.
\eref{107}{1}

\paragraph{Outcome.}
\textbf{Outcome: Exploratory approach with unresolved gap.}

This attempt establishes an elementary independent-point construction over
explicit multiquadratic fields, proves exactly why its points do not descend
as a high-rank rational family, and then obtains an unconditional, exact
rank-at-least-four certificate for one rational curve by a reproducible
same-twist search. These are checks of the proposed mechanism, not a new
rank record or a proof of unbounded rational rank. The missing step is
unbounded certificate rank on one rational curve at a time, or a genuinely
different construction not subject to the same-square-class obstruction.
% END RESEARCH ADDITION 107
\subsection{Sources}
\begin{itemize}\raggedright
\item[\textnormal{[S1]}]\hypertarget{source:107:1}{} Harris B. Daniels and Hannah Goodwillie, On the ranks of elliptic curves with isogenies, arXiv1611.01329v2,11December2024,10pages.
\url{https://arxiv.org/pdf/1611.01329v2}.
{\small\textit{Catalogue formulation source.}}
% BEGIN REFERENCE ADDITION 107
\item[\textnormal{[E1]}]\hypertarget{extra:107:1}{} J. Park, B. Poonen,
J. Voight, and M. Matchett Wood, \emph{A heuristic for boundedness of ranks
of elliptic curves}, arXiv:1602.01431v3 (10 July 2018),
Section 1.1, Remark 1.1.2, and Section 8.2.
\url{https://arxiv.org/html/1602.01431v3}.
{\small\textit{Consulted 27 September 2026; model theorems are distinguished from predictions for actual curves.}}
\item[\textnormal{[E2]}]\hypertarget{extra:107:2}{} A. Dujella,
\emph{History of elliptic curves rank records}, and its
\emph{Rank $\geq31$} data page, crediting Alp\"oge--Howell (2026).
\url{https://web.math.pmf.unizg.hr/~duje/tors/rankhist.html};
\url{https://web.math.pmf.unizg.hr/~duje/tors/rk31.html}.
{\small\textit{Consulted 27 September 2026. The record page and point list were inspected; their complete independence proof was not independently audited.}}
% END REFERENCE ADDITION 107
\end{itemize}

\end{document}
